\section{Routers}\label{sec:methods}

Every router computes the round's tree at the sink from the information in A8--A11 and broadcasts it in the advertisement. All routers in the main comparison build multi-hop trees with full forwarding.

\subsection{Shortest-path baselines}
\emph{Min-energy (MTE) Dijkstra} uses the link cost $e_{\mathrm{tx}}(i,j)+e_{\mathrm{rx}}$. \emph{Battery-weighted Dijkstra} multiplies it by $E_0/E_i$, the sender's inverse residual-battery fraction, and is recomputed every round:
\begin{equation}
w(i,j)=\bigl[e_{\mathrm{tx}}(i,j)+e_{\mathrm{rx}}\bigr]\,E_0/E_i.\label{eq:battery}
\end{equation}
In battery-weighted Dijkstra every node's path is computed independently, so within one round many nodes can select the same relay.

\subsection{Load-balanced shortest-path tree (Shan et al.)}
We implement the top-down construction of~\cite{shan2013lifetime}. Hop layers are built over a radio range $R$. Each new layer is attached to the previous one so that the largest branch load at the sink is minimised, using a bisection on the bound with a max-flow feasibility test and unit load per sensor, as in their model. Within the chosen branch a node takes the candidate parent with the fewest descendants. The tree is rebuilt when a node dies. We do not implement their distributed refinement step. $R$ is tuned on validation deployments ($R=110$~m).

\subsection{Flow augmentation (Chang and Tassiulas)}\label{sec:fa}
Flow augmentation~\cite{chang2004maximum} routes traffic in small increments on shortest-cost paths, updating the residual energy after each increment. We run it as a per-round planner. Over a horizon of $H$ rounds, each source's traffic is pushed in $K$ increments on shortest paths over the cost
\begin{equation}
c(i,j)=e_{\mathrm{tx}}(i,j)\,(E_0/E'_i)^x + e_{\mathrm{rx}}\,(E_0/E'_j)^x,\label{eq:fa}
\end{equation}
where $E'$ is the residual energy predicted after the increments already placed. Each node then forwards to the next hop that carries the largest share of its augmented flow. No solver is used. Validation chose $x=1$, $H=3$ and $K=30$ (Section~\ref{sec:setup}). Because $E'$ changes during the round's own increments, later increments avoid relays loaded by earlier ones. This is a form of within-round coordination.

\subsection{Booked-load Dijkstra}\label{sec:bld}
Booked-load Dijkstra coordinates one round's routes through \emph{load} instead of predicted energy. The sink processes nodes in increasing distance from the sink. For each node without a route it runs Dijkstra over
\begin{equation}
w(j,k)=\bigl[e_{\mathrm{tx}}(j,k)+e_{\mathrm{rx}}\bigr]\,\frac{E_0}{E_j}\Bigl(1+\kappa\,\frac{A_j}{\bar r}\Bigr),\label{eq:bld}
\end{equation}
where $A_j$ is the traffic already booked through node $j$ in this round's construction and $\bar r$ is the mean node rate. It then fixes the path and books the source's rate on every node of it. A node that already received a route as a relay keeps its next hop, and its own rate is booked along that fixed path when its turn comes. Rates $r_i$ are the sink's count of packets received from $i$ over the last 24 rounds, so no extra bits are needed. With $\kappa=0$ the method is battery-weighted Dijkstra. Validation chose $\kappa=0.5$, and Section~\ref{sec:kappa} sweeps $\kappa$ on fresh deployments. Algorithm~\ref{alg:bld} gives one round.

\begin{algorithm}[t]
\caption{Booked-load Dijkstra: one round at the sink}\label{alg:bld}
\small
\begin{algorithmic}[1]
\Require alive nodes $V$, batteries $E_j$, observed rates $r_j$, link energies, weight $\kappa$
\Ensure parent function $p$
\State $A_j\gets0$, $p(j)\gets\bot$ for all $j$; $\bar r\gets$ mean of $r_j$
\For{$i\in V$ in increasing distance to the sink}
  \If{$p(i)\ne\bot$}
    \State $u\gets i$; \textbf{while} $u\ne s$: $A_u\gets A_u+r_i$, $u\gets p(u)$ \Comment{book own traffic}
    \State \textbf{continue}
  \EndIf
  \State arc weights from~\eqref{eq:bld}; for $j$ with $p(j)\ne\bot$ keep only $j\to p(j)$
  \State $P\gets$ Dijkstra least-cost path from $i$ to $s$
  \For{each $u$ on $P$ except $s$}
    \If{$p(u)=\bot$} $p(u)\gets$ successor of $u$ on $P$ \EndIf
    \State $A_u\gets A_u+r_i$
  \EndFor
\EndFor
\State \Return $p$
\end{algorithmic}
\end{algorithm}

\subsection{Ablations of booked-load Dijkstra}\label{sec:lest}
\emph{Load only} removes the battery factor from~\eqref{eq:bld}. \emph{Tiered load} replaces the sink-observed rates with load reported in four tiers (boundaries 0.15, 0.40, 0.75 of twice the largest historical rate, with a 0.05 hysteresis band). This is the Lightweight State Table design of our earlier work, and it is charged as in Section~\ref{sec:packet}. In a centralised sink the observed rates are free, so the tiered variant cannot gain. It is reported to quantify what quantised sharing would cost in a setting where nodes, not the sink, need the load.

\subsection{Upper reference: LP planner}\label{sec:lp-planner}
The sink re-solves the LP~\eqref{eq:lp-flow}--\eqref{eq:lp-energy} every round from the current batteries. Each node generates its observed mean rate, and a small total-energy term $\mu=10^{-3}$ favours low-energy flows. Every node then sends to the next hop with the largest LP share; credits are reset at each solve, so this is a deterministic rounding of the flow, not proportional rotation. It uses the same information as the routers above and is the strongest realisable router we found. It needs an LP solver at the sink.

\begin{widetab}
\centering
\caption{What each router needs ($n=100$). Times are single-thread wall-clock per round, mean over the main bench.}\label{tab:characteristics}
\footnotesize
\begin{adjustbox}{max width=\linewidth}
\begin{tabular}{@{}L{0.25\textwidth}L{0.24\textwidth}L{0.22\textwidth}L{0.12\textwidth}r@{}}
\toprule
Router & Information at the sink & Work per round & Extra bits & ms/round\\
\midrule
\inputtab{body/rev_char.tex}
\bottomrule
\end{tabular}
\end{adjustbox}
\end{widetab}
